% 359_KI_Grenzen_Frequenzanalyse_De_ch.tex
% Johann Pascher, 8. September 2026
% KI-basierte Mustererkennung und algebraische Grenzen der Frequenzanalyse
\providecommand{\BM}{\textbf{[B]}}
\providecommand{\KM}{\textbf{[K]}}
\providecommand{\SM}{\textbf{[S]}}
\providecommand{\XM}{\textbf{[X]}}
\providecommand{\QM}{\textbf{[Q]}}
\providecommand{\EM}{\textbf{[E]}}
\providecommand{\LM}{\textbf{[L]}}

\phantomsection
\addcontentsline{toc}{chapter}{%
  KI-basierte Mustererkennung und algebraische Grenzen der Frequenzanalyse}

\section*{Statuszeichen}
\begin{center}
\begin{tabular}{@{}lp{10.5cm}@{}}
\textbf{[B]} & \emph{Belegt} --- algebraisch bewiesen. \\[3pt]
\textbf{[K]} & \emph{Konfirmiert} --- numerisch geprüft. \\[3pt]
\textbf{[S]} & \emph{Spekulativ} --- offen. \\[3pt]
\textbf{[E]} & \emph{Etabliert} --- Lehrbuch- oder Konsensuswissen. \\[3pt]
\textbf{[L]} & \emph{Literaturbeleg} --- externer Primärnachweis. \\[3pt]
\textbf{[Q]} & \emph{Quelle} --- externer Primärwert. \\
\end{tabular}
\end{center}

\section*{Abstract}

KI-basierte Mustererkennung stößt bei der Frequenzanalyse an dieselben
algebraischen Grenzen wie klassische Methoden --- aber aus einem anderen
Grund und ohne es zu wissen. Dieses Dokument benennt die drei Grenzen
präzise, belegt sie mit aktueller Literatur, und skizziert, was eine
Galois-informierte KI-Architektur leisten könnte.

\textbf{Ergebnisse:}
\begin{enumerate}
\item Satz~A \BM: Die Standard-HRV-Frequenzbänder (Task Force 1996) haben
  Bandgrenzen ausschließlich mit 5-Limit-Primfaktoren $\{2,3,5\}$ und
  liegen damit im Galois-Raster von Dok.~358 Satz~B. Sie sind mit dem
  Raster verträglich, aber nicht durch es festgelegt: Gerundete
  Dezimalwerte liegen fast immer im 5-Limit (Dok.~360).
\item Satz~B \BM: Die Auflösungsgrenze $100\xi\approx1{,}33\,\%$ (Dok.~343)
  hat die Farey-Höhe $Q_c\approx11{,}1$. Eine Klassifikation mit dieser
  Toleranz verlangt eine Frequenzauflösung $\Delta f/f\le100\xi$.
\item Satz~C \BM: Gruppenäquivariante Netze (G-CNNs) für kontinuierliche
  Gruppen (E(3), SO(3)) reduzieren Parameteranzahl und verbessern
  Generalisierung nachweislich. Für endliche Galois-Gruppen GF$(3^k)$
  existiert diese Architektur noch nicht --- das ist die identifizierte Lücke.
\item Satz~D \BM: Ob das Zulässigkeitskriterium aus Dok.~358 Satz~B
  lernbar ist, entscheidet die Hypothesenklasse, nicht die Datenmenge.
  Ohne Strukturannahme verlangt die PAC-Schranke rund 84\,000 markierte
  Beispiele; mit der Strukturannahme \enquote{Primfaktoren aus einer
  festen Menge} genügen höchstens 407, und aus positiven Beispielen allein
  wird die Primzahlmenge $\{2,3,5,7,11,13\}$ schon nach 30 Beispielen
  sicher erkannt.
\item Architekturskizze \SM: Ein Galois-informiertes neuronales Netz
  verwendet GF$(3^k)$-Orbits als hard constraint in der Eingangsschicht.
  Das reduziert den Parameterraum und macht das Zulässigkeitskriterium
  nicht lernbar, sondern erzwungen.
\end{enumerate}

Prüfskript: \texttt{pruef\_359\_ki\_grenzen.py} (18/18).

% ============================================================
\section{Ausgangslage: was KI heute tut}
% ============================================================

\subsection*{Physics-Informed Neural Networks (PINNs)}

PINNs verwenden bekannte Differentialgleichungen als Hard-Constraint
in der Verlustfunktion \LM\,\cite{Raissi2024}. Sie kennen keine
algebraische Gruppenstruktur. Ein PINN kann eine Schwingungsgleichung
einbauen, nicht aber die Bedingung ``Primfaktoren nur in
$\{2,3,5,7,11,13\}$''.

\subsection*{Gruppenäquivariante Netze (G-CNNs)}

G-CNNs bauen Symmetriegruppen direkt in die Architektur ein, so dass
der Output äquivariant unter Gruppenoperationen ist \LM\,\cite{Finzi2022}.
Für kontinuierliche Gruppen wie E$(3)$ (Rotation, Translation) und
SO$(3)$ sind solche Architekturen seit 2021 etabliert und zeigen
messbare Verbesserungen in Molekülmodellierung und Materialwissenschaften.

Für \emph{endliche} Körper GF$(3^k)$ fehlt eine vergleichbare Architektur
vollständig. Das ist die Lücke, die dieses Dokument identifiziert.

\subsection*{HRV-Analyse mit KI}

Multi-resolution-HRV-Analyse mit Machine Learning zeigt, dass verschiedene
Auflösungsstufen unterschiedliche Merkmale liefern \LM\,\cite{Gupta2024}. Aktuelle
KI-HRV-Systeme verwenden die Task-Force-Bänder VLF/LF/HF als feste
Eingabegröße, ohne zu wissen, warum diese Grenzen dort liegen
\LM\,\cite{Villanueva2026}.

% ============================================================
\section{Satz A: HRV-Bandgrenzen im Galois-Raster}
% ============================================================

\begin{quote}
\textbf{Satz A.} \BM\ Die Standard-HRV-Frequenzbänder \LM\,\cite{TaskForce1996}
haben folgende Grenzen als einfachste rationale Näherungen:
\begin{align}
f_\text{VLF,max} &= \tfrac{1}{25}\,\text{Hz} = 0{,}04\,\text{Hz}
  && \text{Primfaktoren: } \{5\} \\
f_\text{LF,max}  &= \tfrac{3}{20}\,\text{Hz} = 0{,}15\,\text{Hz}
  && \text{Primfaktoren: } \{2,3,5\} \\
f_\text{HF,max}  &= \tfrac{2}{5}\,\text{Hz} = 0{,}40\,\text{Hz}
  && \text{Primfaktoren: } \{2,5\}
\end{align}
Alle Primfaktoren liegen in $\{2,3,5\}$ --- dem 5-Limit-Raster und
gleichzeitig im Galois-Raster $\{2,3,5,7,11,13\}$ von Dok.~358 Satz~B.
Die Bandverhältnisse $f_\text{LF}/f_\text{VLF} = 15/4$ und
$f_\text{HF}/f_\text{LF} = 8/3$ liegen ebenfalls im 5-Limit.
\end{quote}

\emph{Beweis.} Direkte Primfaktorisierung (Prüfskript A1--A5). \qed

\textbf{Konsequenz für KI:} Aktuelle HRV-KI-Modelle verwenden diese
Bandgrenzen als freie Hyperparameter, die aus Daten gelernt oder manuell
gesetzt werden \LM\,\cite{Villanueva2026}. Ein Galois-informiertes Netz
könnte sie als feste Konstanten übernehmen. Das ist eine Setzung, keine
algebraische Notwendigkeit: 8 von 9 Zahlen mit einer gültigen Ziffer und
ein Drittel aller zweistelligen Dezimalwerte liegen im 5-Limit, und in
rad/s sind die Grenzen irrational (Dok.~360). Als ``inductive bias'' im
Sinne der Lerntheorie \LM\,\cite{ILCRBlog2023} ist die Übernahme
zulässig, aber nicht algebraisch begründet.

% ============================================================
\section{Satz B: Auflösungsgrenze und Farey-Höhe}
% ============================================================

\begin{quote}
\textbf{Satz B.} \BM\ Die FFGFT-Auflösungsgrenze $100\xi\approx1{,}33\,\%$
(Dok.~343 Satz~G$'$) hat Farey-Nennerhöhe
$Q_c = \pi/\sqrt{6\cdot100\xi} \approx 11{,}1$.
Frequenzverhältnisse mit Nennern $>Q_c$ sind im Galois-Raster nicht
trennbar. Die HRV-Bandgrenzen haben Nenner 5, 20, 25 --- alle im
trennbaren Bereich. Eine Klassifikation mit der Toleranz $100\xi$ ist
nur sinnvoll, wenn die Frequenzauflösung $\Delta f/f\le100\xi$ erfüllt.
\end{quote}

\textbf{Praktische Folgerung \KM:} Mit $\Delta f=1/T_w$ für die
Fensterlänge $T_w$ (bei Welch die Segmentlänge) folgt $T_w\ge1/(100\xi\,f)$,
also 12{,}5~min bei 0{,}10~Hz und 15~min bei 0{,}083~Hz (Dok.~360).
Kürzere Fenster können Rasterelemente nicht unterscheiden. Ob die
Schwelle in KI-HRV-Analysen empirisch auftritt, ist nicht belegt \SM.

% ============================================================
\section{Satz C: Was gruppenäquivariante Netze leisten --- und wo die Lücke ist}
% ============================================================

\begin{quote}
\textbf{Satz C.} \BM\ Ein G-equivariantes Netz mit Gruppe $G$ hat in
einer vollständig verbundenen Schicht $\mathbb{R}^{|G|}\to\mathbb{R}^{|G|}$
nur $|G|$ freie Parameter statt $|G|^2$, weil die Gewichtsmatrix
durch die Gruppenstruktur vollständig bestimmt ist \LM\,\cite{Finzi2022}.
Für $G = (\mathbb{Z}/13)^*$ mit $|G|=12$ bedeutet das: Faktor~12
Parameterreduktion gegenüber einem ungruppierten Netz.
\end{quote}

\emph{Beweis.} Die Gruppe $(\mathbb{Z}/13)^*$ ist zyklisch der Ordnung~12
mit Erzeugern $\{2,6,7,11\}$ (Prüfskript C2). Ein $G$-equivariantes Netz
hat Gewichtsmatrix der Form $(W_{ij}) = w_{j-i \bmod 12}$ --- ein
zirkulantes Netz mit $|G|=12$ Parametern (Prüfskript C3). \qed

\textbf{Was es für E$(3)$ und SO$(3)$ gibt \LM\,\cite{Finzi2022}:}
Equivariante Netze für kontinuierliche Gruppen (Rotationen, Translationen)
sind in der Molekülmodellierung und Materialwissenschaft Standard.
Sie verwenden sphärische Harmonische und Clebsch-Gordan-Produkte als
Bausteine.

\textbf{Was fehlt:} Eine analoge Architektur für GF$(3^k)^*$ --- eine
endliche, nicht-kontinuierliche Gruppe. Die Herausforderung: GF$(3^k)^*$
hat keine natürliche Einbettung in einen euklidischen Raum, auf dem
Standard-KI-Schichten operieren. Die Orbits der Frobenius-Wirkung
($x\mapsto x^3$) müssen als diskrete Symmetrien in die Architektur
eingebaut werden, nicht als kontinuierliche Transformationen.

\textbf{Architekturskizze \SM:} (Abschnitt~\ref{sec:arch})

% ============================================================
\section{Satz D: Lernbarkeit hängt an der Hypothesenklasse}
% ============================================================

\begin{quote}
\textbf{Satz D.} \BM\ Auf dem beschränkten Raster der gekürzten Brüche
$p/q$ mit $p,q\le100$ ist das Zulässigkeitskriterium aus Dok.~358 Satz~B
(Primfaktoren nur in $\{2,3,5,7,11,13\}$) lernbar. Wie viele Beispiele
dafür nötig sind, entscheidet die Hypothesenklasse, also das Vorwissen,
das ein lernendes System mitbringt, und nicht die Datenmenge allein.
\end{quote}

\emph{Begründung.} Das Raster umfasst 6087 gekürzte Brüche, davon sind
1453 zulässig (Prüfskripte D1, D3). Hundert Beispiele decken weniger als
$7\,\%$ der zulässigen Verhältnisse ab (D2). Für eine endliche
Hypothesenklasse $\mathcal H$ genügen im realisierbaren Fall
$m\ge\bigl(\ln|\mathcal H|+\ln(1/\delta)\bigr)/\varepsilon$ markierte
Beispiele, um mit Wahrscheinlichkeit $1-\delta$ einen Fehler unter
$\varepsilon$ zu erreichen \LM\,\cite{ShalevShwartz2014}. Mit
$\varepsilon=\delta=5\,\%$ ergibt das zwei sehr verschiedene Fälle.

Ohne Strukturannahme, wenn also jede Teilmenge der 6087 Brüche als
Kriterium in Frage kommt, ist $|\mathcal H|=2^{6087}$ und die
VC-Dimension gleich 6087. Die Schranke verlangt rund 84\,000 markierte
Beispiele (D4), und aus positiven Beispielen allein ist das Kriterium
nicht zu gewinnen, weil jedes ungesehene Verhältnis ebenso gut zulässig
sein könnte. Das ist die Lage eines generischen Netzes ohne
eingebautes Vorwissen.

Mit der Strukturannahme \enquote{zulässig sind genau die Brüche, deren
Primfaktoren in einer festen Menge $P$ liegen}, wobei $P$ eine Teilmenge
der 25 Primzahlen bis 100 ist, gilt $|\mathcal H|=2^{25}$, und es
genügen höchstens 407 Beispiele (D5). In dieser Klasse lernt sogar ein
Verfahren, das nur positive Beispiele sieht: Es nimmt als $P$ die
kleinste Primzahlmenge, die alle gesehenen Beispiele erklärt. Bei
zufällig gezogenen zulässigen Verhältnissen erkennt es
$\{2,3,5,7,11,13\}$ nach 10 Beispielen in $94{,}6\,\%$ und nach
30 Beispielen in allen von 1000 Durchläufen (D6). Ein solches System
schließt $17/11$ aus, ohne je ein Negativbeispiel gesehen zu haben, weil
die 17 in keinem Beispiel auftritt.

Der Satz von Gold \LM\,\cite{Gold1967} betrifft die Identifikation einer
Sprache im Grenzwert aus unendlichen Folgen positiver Beispiele; er
greift erst, wenn die Hypothesenklasse unbeschränkt ist, und trägt für
das beschränkte Raster keine Nicht-Lernbarkeit.

\textbf{Was das heißt:} Das Wissen, dass nur die Primzahlen bis 13
vorkommen, steckt nicht in der Menge der Trainingsdaten, sondern in der
Form der Hypothesen. Ein Netz, dem diese Form fehlt, braucht um mehr als
zwei Größenordnungen mehr markierte Beispiele, und mit Messrauschen
oberhalb der Auflösungsgrenze aus Satz~B gewinnt es die fehlende
Information gar nicht.

\textbf{Konsequenz:} Das Zulässigkeitskriterium gehört als Strukturvorgabe
in die Architektur, nicht als nachträgliche Regularisierung.
Punkt-Regularisierung ($\ell_1$, $\ell_2$) bevorzugt kleine Gewichte,
aber keine Primfaktormenge, und erzwingt deshalb keine algebraische
Unzulässigkeit.

% ============================================================
\section{Architekturskizze: Galois-informiertes neuronales Netz}
\label{sec:arch}
% ============================================================

Die folgende Skizze ist spekulativ \SM\ und soll die Richtung zeigen,
nicht eine fertige Architektur beschreiben.

\subsection*{Eingangsschicht: Galois-Projektion}

Die Eingangsschicht bildet gemessene Frequenzverhältnisse $f_i/f_j$
auf die nächstliegende rationale Zahl im Galois-Raster ab:
\[
  \pi_G(x) = \underset{p/q \in \mathcal{R}_G}{\arg\min}\,|x - p/q|,
  \quad \mathcal{R}_G = \{p/q : \text{Primfaktoren} \subseteq \{2,3,5,7,11,13\}\}
\]
Das ist keine Lernoperation --- es ist ein algebraischer Vorverarbeitungsschritt.
Verhältnisse außerhalb des Rasters werden auf ihre nächste Raster-Approximation
projiziert und mit einem Distanzflag versehen.

\subsection*{Mittlere Schichten: $(\mathbb{Z}/13)^*$-Equivarianz}

Die mittleren Schichten sind $(\mathbb{Z}/13)^*$-equivariant: zirkulante
Matrizen mit 12 statt 144 Parametern pro Schicht. Das entspricht genau
dem Prinzip der G-CNNs für die Gruppe $G = (\mathbb{Z}/13)^*$.
Die Frobenius-Wirkung $x\mapsto x^3$ ist die natürliche Symmetrie
(Dok.~336, 341) und wird als Gewichtsteilung eingebaut.

\subsection*{Auflösungsfilter: $100\xi$-Schwelle}

Vor der Galois-Projektion: ein Tiefpassfilter mit Grenzfrequenz
$\Delta f/f = 100\xi \approx 1{,}33\,\%$. Alles unterhalb dieser Schwelle
wird als Einbettungsrauschen behandelt und nicht klassifiziert.
Die Schwelle ist eine Setzung der FFGFT; eine empirische Entsprechung in
HRV-Daten ist nicht belegt (Dok.~360).

\subsection*{Stabilitätsbewertung: Arnold-Zungenbreite}

Der Output des Netzes enthält nicht nur eine Klassifikation, sondern
auch eine Stabilitätsbewertung: die algebraisch berechnete
Arnold-Zungenbreite (Dok.~328) des klassifizierten Verhältnisses.
Das gibt dem Anwender sofort eine Unsicherheitsschätzung --- nicht
aus dem Netz gelernt, sondern aus der Algebra berechnet.

\subsection*{Was das leistet gegenüber Standard-KI}

\begin{center}
\begin{tabular}{@{}lll@{}}
\toprule
& Standard-KI & Galois-informiert \\
\midrule
Zulässigkeitskriterium & gelernt (unvollständig) & erzwungen (vollständig) \\
Auflösungsgrenze & nicht vorgesehen & algebraisch gesetzt \\
Stabilitätsvorhersage & nur für gesehene Fälle & für alle Fälle (Arnold-Zunge) \\
Parameteranzahl & $|G|^2 = 144$ & $|G| = 12$ \\
Interpretierbarkeit & beschränkt & vollständig \\
\bottomrule
\end{tabular}
\end{center}

% ============================================================
\section{Anwendungsfall HRV}
% ============================================================

Für die Herzratenvariabilitätsanalyse ergeben sich konkrete Konsequenzen:

\begin{enumerate}
\item \textbf{Bandgrenzen als feste Konstanten:} VLF/LF/HF-Grenzen
  werden nicht aus Daten gelernt (Satz~A); sie sind mit dem Raster
  verträglich, aber eine Setzung. Der Suchraum des Netzes reduziert sich
  entsprechend.
\item \textbf{Erforderliche Auflösung:} $\Delta f/f \le 1{,}33\,\%$
  (Satz~B), also $\approx0{,}0013$\,Hz bei 0{,}10~Hz und eine
  Fensterlänge von mindestens 12{,}5~min.
\item \textbf{Neue Biomarker:} Statt LF/HF-Ratio (ein einzelner Zahlenwert)
  könnte ein Galois-informiertes Netz den \emph{Orbit-Typ} des dominanten
  Frequenzverhältnisses ausgeben --- ein algebraischer Biomarker, der
  angibt, in welcher Galois-Schicht ($k=3,4,5,6$) das Herz aktuell
  schwingt. Das ist neu und hat keine Entsprechung in der Literatur \SM.
\item \textbf{Stabilitätsvorhersage \SM:} Ob ein aktuelles HRV-Muster
  robust ist (breite Arnold-Zunge) oder fragil (schmale Zunge). Die
  Prämisse, dass HRV-Frequenzverhältnisse auf dem Raster einrasten, ist
  ungeprüft; die Daten in Dok.~360 stützen sie nicht, ein prospektiver
  Test ist dort beschrieben.
\end{enumerate}

% ============================================================
\section{Abgrenzung}
% ============================================================

\begin{itemize}
\item Dieses Dokument implementiert keine KI-Architektur und macht keine
  klinischen Aussagen.
\item Die Architekturskizze in §5 ist spekulativ \SM\ --- sie skizziert
  eine Richtung, kein fertiges System.
\item Der algebraische Orbit-Biomarker in §6 ist eine Hypothese; seine
  physiologische Bedeutung ist offen. Die explorative Prüfung an HRV-Daten
  (Dok.~360) ergab einen Negativbefund.
\item Satz~D vergleicht Stichprobenschranken auf dem beschränkten
  Raster $p,q\le100$; über einem unbeschränkten Raster mit beliebig
  großen Primzahlen ist die Hypothesenklasse unendlich, und die Aussagen
  sind neu zu prüfen.
\end{itemize}

% ============================================================
\section{Zusammenfassung}
% ============================================================

\begin{center}
\begin{tabular}{@{}lp{9cm}l@{}}
\toprule
Satz & Aussage & Status \\
\midrule
A & HRV-Bandgrenzen im 5-Limit-Raster: verträglich, nicht festgelegt & \BM \\
B & Auflösungsgrenze $100\xi\approx1{,}33\,\%$, $Q_c\approx11{,}1$; Klassifikation verlangt $\Delta f/f\le100\xi$ & \BM \\
C & G-CNNs für GF$(3^k)$: Parameterreduktion 12x; Architektur fehlt & \BM/\SM \\
D & Lernbarkeit des Zulässigkeitskriteriums: ohne Strukturannahme $\sim$84\,000, mit Primfaktor-Struktur $\le$407 Beispiele & \BM \\
Arch. & Galois-informiertes Netz: Skizze & \SM \\
\bottomrule
\end{tabular}
\end{center}

\begin{thebibliography}{9}
\bibitem{Raissi2024}
  Raissi, M., Perdikaris, P., Ahmadi, N., Karniadakis, G.\,E. (2024).
  \emph{Physics-Informed Neural Networks and Extensions.}
  arXiv:2408.16806.
\bibitem{Finzi2022}
  Finzi, M. (2023).
  \emph{Understanding and Incorporating Mathematical Inductive Biases
  in Neural Networks.} PhD Thesis, NYU.
  \url{https://cs.nyu.edu/media/publications/Marc_Finzi_Thesis__12_.pdf}
\bibitem{Gupta2024}
  Gupta, K., Sinha, G.\,R., Bhat, R., Saoji, A.\,A., \& Manjunath, N.\,K.\ (2024).
  Multi-resolution assessment of heart rate variability signals during
  yogic and normal breathing using machine learning modules.
  \emph{Applied Soft Computing}, 164, 112020.
  DOI: 10.1016/j.asoc.2024.112020.
\bibitem{Villanueva2026}
  Zimatore, G.\ et al.\ (2026).
  HRV in Stress Monitoring by AI: A Scoping Review.
  \emph{MDPI Applied Sciences}, 16(1):23.
  \url{https://www.mdpi.com/2076-3417/16/1/23}
\bibitem{TaskForce1996}
  Task Force of the European Society of Cardiology and the North American
  Society of Pacing and Electrophysiology (1996).
  Heart rate variability: standards of measurement, physiological
  interpretation and clinical use.
  \emph{European Heart Journal}, 17, 354--381.
\bibitem{Gold1967}
  Gold, E.\,M. (1967).
  Language identification in the limit.
  \emph{Information and Control}, 10(5), 447--474.
\bibitem{ShalevShwartz2014}
  Shalev-Shwartz, S., \& Ben-David, S. (2014).
  \emph{Understanding Machine Learning: From Theory to Algorithms.}
  Cambridge University Press.
\bibitem{ILCRBlog2023}
  ICLR Blog (2023).
  How does the inductive bias influence the generalization capability
  of neural networks?
  \url{https://iclr-blogposts.github.io/2023/blog/2023/how-does-the-inductive-bias-influence-the-generalization-capability-of-neural-networks}
\bibitem{Lowet2022}
  Lowet, E.\ et al.\ (2022).
  Tuning Neural Synchronization: The Role of Variable Oscillation
  Frequencies in Neural Circuits.
  \emph{Frontiers in Systems Neuroscience}.
  \url{https://doi.org/10.3389/fnsys.2022.908665}
\end{thebibliography}

\section*{Prüfskripte}

\texttt{2/python/Dok359\_Skripte/pruef\_359\_ki\_grenzen.py} ---
18 Assertions: A1--A5 (HRV-Raster), B1--B4 (Auflösung), C1--C3
(Gruppenäquivarianz), D1--D6 (Raster, Abdeckung, PAC-Schranken,
Lernen aus positiven Beispielen).

\section*{Vermerk zur Verwendung von KI-Assistenz}

Bei der Erstellung dieses Dokuments wurde KI-Assistenz (Claude, Anthropic)
für Formulierung, LaTeX-Satz, Literaturrecherche und die Erstellung des
Prüfskripts verwendet. Die inhaltlichen Aussagen, die Statuszuordnung
und die Verantwortung für den Text liegen beim Autor.
